% Part: proof-theory
% Chapter: normalization
% Section: introduction

\documentclass[../../../include/open-logic-section]{subfiles}

\begin{document}

\olfileid{pt}{nor}{seg}

\olsection{खंड आणि कट}

!!^a{introduction} नियमाच्या
नंतर !!a{elimination} नियम
येणे हा !!a{proof} मधला
वळसा असतो; !!a{proof}
सामान्यरूप करणे म्हणजे
असे वळसे काढणे.
पण~\Elim{\lor} आणि
\Elim{\lexists} मुळे
“काढायचा वळसा” याची
व्याख्याही थोडी
गुंतागुंतीची होते.
अडचण अशी: वळशात
!!{elimination} नियम
!!{introduction} नियमाच्या
\emph{लगेच} नंतर येईलच
असे नाही. त्यांच्यामध्ये
\Elim{\lor} किंवा
\Elim{\lexists} असू
शकतो; उदाहरणार्थ:
\[
\AxiomC{}
\DeduceC{$\lexists[x][!C(x)]$}
\AxiomC{$\Discharge{!A}{x}$}
\DeduceC{$!B$}
\DischargeRule{\Intro{\lif}}{x}
\UnaryInfC{$!A \lif !B$}
\DischargeRule{\Elim{\lexists}}{y}
\BinaryInfC{$!A \lif !B$}
\AxiomC{}
\DeduceC{$!A$}
\RightLabel{\Elim{\lif}}
\BinaryInfC{$!B$}
\DisplayProof
\]
इतकेच नव्हे,~\Intro{\lif}
आणि~\Elim{\lif} यांच्या
मध्ये कितीही
\Elim{\lor} किंवा
\Elim{\lexists} अनुमाने
असू शकतात; उदाहरणार्थ:
\[
\AxiomC{}
\DeduceC{$\lexists[x][!C(x)]$}
\AxiomC{}
\DeduceC{$!D \lor !E$}
\AxiomC{$\Discharge{!A}{x}$}
\DeduceC{$!B$}
\DischargeRule{\Intro{\lif}}{x}
\UnaryInfC{$!A \lif !B$}
\AxiomC{$\Discharge{!A}{x}$}
\DeduceC{$!B$}
\DischargeRule{\Intro{\lif}}{x}
\UnaryInfC{$!A \lif !B$}
\RightLabel{\Elim{\lor}}
\TrinaryInfC{$!A \lif !B$}
\DischargeRule{\Elim{\lexists}}{y}
\BinaryInfC{$!A \lif !B$}
\AxiomC{}
\DeduceC{$!A$}
\RightLabel{\Elim{\lif}}
\BinaryInfC{$!B$}
\DisplayProof
\]
या उदाहरणात तोच
\Elim{\lif} एकाहून
अधिक वळशांत भाग
घेतो: त्याच्या आधी
एकाहून अधिक~\Intro{\lif}
आहेत. अशा सर्व
शक्यता सामावण्यासाठी
~\Log{N1i}-सिद्धतेतील
!!{formula} आस्थितींच्या
विशिष्ट अनुक्रमांनी
वळसे परिभाषित करतो;
त्यांना \emph{खंड}
म्हणतात. येथे त्या
~$!A \lif !B$
च्या आस्थितींच्या
अनुक्रमिका आहेत:
$\Elim{\lor}$ किंवा
$\Elim{\lexists}$
च्या गौण आधारविधानानंतर
त्याच अनुमानाचा
निष्कर्ष येतो.
या उदाहरणात असे
दोन खंड आहेत;
प्रत्येकात~$!A \lif !B$
च्या तीन आस्थिती
आहेत. आता औपचारिक
व्याख्या पाहू.

\begin{defn}
~\Log{N1i} मधील
!!{proof}~$\delta$
मध्ये एकाच !!{formula}~$!A$
च्या $!A_1$, \dots,
$!A_n$ या~$\delta$
मधील आस्थितींचा
अनुक्रम पुढील अटी
पूर्ण करत असेल,
तर तो \emph{खंड} आहे:
\begin{enumerate}
  \item $!A_1$ हा
  \Elim{\lor} किंवा
  \Elim{\lexists}
  चा निष्कर्ष नसतो;
  \item $!A_n$ हे
  \Elim{\lor} किंवा
  \Elim{\lexists}
  चे गौण आधारविधान नसते;
  \item $i = 1$, \dots,~$n-1$
  साठी~$!A_i$ हे
  \Elim{\lor} किंवा
  \Elim{\lexists}
  चे गौण आधारविधान
  असते, आणि~$!A_{i+1}$
  हा त्याच अनुमानाचा
  निष्कर्ष असतो.
\end{enumerate}
त्या खंडाची
\emph{लांबी}~$n$ आहे.
\end{defn}

प्रत्येक खंड काढून
टाकण्याजोगा वळसा
असतोच असे नाही.
अशा वळशांना
\emph{कट} म्हणतात.

\begin{defn}
एखाद्या खंडातील
$!A_n$ हे
!!a{elimination} अनुमानाचे
प्रमुख आधारविधान असेल,
आणि पुढीलपैकी एक
अट पूर्ण होत असेल,
तर तो \emph{कट} आहे:
$n=1$ असून
$!A_1 = !A_n$ हे
!!a{introduction} अनुमानाचा
किंवा~\FalseInt चा
निष्कर्ष असेल;
किंवा~$n>1$ असेल.
\end{defn}

कट-खंडाची सुरुवात
!!a{introduction} नियमाच्या
निष्कर्षाने व्हावीच
लागते असे नाही.
तो !!a{elimination}
नियमाच्या प्रमुख
आधारविधानात संपणे
हीच आवश्यक अट आहे.
मात्र लांबी~$1$
असलेला खंड कट
म्हणून मोजायचा असेल,
तर तो !!{introduction}
नियमाचा किंवा~\FalseInt
चा निष्कर्ष असणेही
आवश्यक आहे.

\begin{defn}
कटाची \emph{कोटी}
ही~$\depth{!A}$
असते. !!a{proof}~$\delta$
ची \emph{कट-कोटी}
$\cutrank{\delta}$
म्हणजे~$\delta$
मधील कटांच्या
कोटींपैकी महत्तम कोटी.
कटाची कोटी
$\cutrank{\delta}$
इतकी असेल, म्हणजे
ती महत्तम असेल,
तर तो~$\delta$
मधील \emph{महत्तम} कट.
~$\delta$ ची
\emph{कट-लांबी}
म्हणजे तिच्यातील
सर्व महत्तम कटांच्या
लांबींची बेरीज.
\end{defn}

\end{document}
